\documentclass[11pt]{article}
\usepackage{fiche}
\matiere{Support Vecteur Machine}
\chapitre{M2}
\source{SVM.pdf}
%% statut: propre

\begin{document}
%% page: SVM.pdf (p. 2)
\titre{\couleur{rose}{\entoure{2} \souligne{Support Vector Machine}}}

\begin{itemize}
\item Rappel : C'est quoi le couple $(w, b)$ optimal ?

\couleur{rouge}{$\hookrightarrow$ Celui qui optimise le critère de la marge optimale}

\couleur{rouge}{\souligne{\# On veut un hyperplan le + loin des données possible}}
\end{itemize}

\dessin[5,15,22,24]{nuage de points (croix et ronds) séparés par un hyperplan hachuré, distance $d$ entre un point et l'hyperplan}

\couleur{rouge}{\} $2 \cdot d$ = \souligne{marge}}

\begin{itemize}
\item Le vecteur support, c'est les observations qui sont sur la marge
\begin{encadre}
\[ SV = \left\{ x \in \mathbb{R}^p : d(H, x) = \frac{\text{marge}}{2} \right\} \]
\end{encadre}

\item La marge est donc défini comme étant :
\begin{encadre}
\[ \text{Marge} = M(w, b) = 2 \cdot d(x_{SV}, H) = 2 \cdot \frac{\overbrace{|h(x_{SV})|}^{\couleur{rouge}{=1}}}{\|w\|} = \entoure{$\dfrac{2}{\|w\|}$} \]
\end{encadre}

\couleur{vert}{$d = \text{Marge}/2 = \dfrac{1}{\|w\|}$}

\item Grâce à ça, on peut maintenant classifier notre échantillon sur la variable cible
\begin{encadre}
\[ y_i = \left\{ \begin{array}{ll} 1 & \text{si } \dfrac{h(x_i)}{\|w\|} \geq m \\[2mm] -1 & \text{si } \dfrac{h(x_i)}{\|w\|} < -m \end{array} \right. \]
\end{encadre}
\end{itemize}

\marge{\couleur{rose}{Rappel} \\ $\|x\| = \sqrt{\sum_{i=1}^{n} x_i^2}$ \\ $\langle x, y \rangle = \sum_{i=1}^{n} x_i y_i$ \\ $\|x\| = \sqrt{\langle x, x \rangle}$}

\begin{itemize}
\item On définit le vecteur normal (normalisé) $\tilde{w}$, la droite orthogonal à l'hyperplan séparateur $H$
\begin{encadre}
\[ \tilde{w} = \frac{w}{\|w\|}, \quad \|\tilde{w}\| = 1 \]
\end{encadre}

\couleur{vert}{$\tilde{w}$ c'est le vecteur normalisé qui donne une distance et qui définit de quel côté de l'hyperplan on \souligne{est}}

\item On peut écrire la formule de $y_i$ selon le vecteur normalisé :

$\langle \tilde{w}, x_i \rangle$ = écart/distance entre le vecteur normal et $x_i$
\begin{encadre}
\begin{align*}
\langle \tilde{w}, x_i \rangle + \tilde{b} \geq 0 &\Rightarrow y_i = 1 \\
\langle \tilde{w}, x_i \rangle \geq \tilde{c} &\Rightarrow y_i = 1
\end{align*}
\end{encadre}

\couleur{rouge}{$\hookrightarrow$ Si on fait le point $x_i$ sur le vecteur $\tilde{w}$ et qui dépasse $\tilde{c}$ $\Rightarrow$ il est au-dessus de l'hyperplan $\Rightarrow y_i = 1$}
\end{itemize}

\marge{\couleur{rose}{\# $\tilde{b} \neq b$ \\ $\tilde{b} = -\tilde{c}$}}

\dessin[69,54,93,62]{hyperplan séparateur, vecteur $\tilde{w}$ partant de l'origine, point A (bleu) et point B (rouge) projetés sur l'hyperplan}

\couleur{rouge}{$d(\tilde{w}, A) = |\langle \tilde{w}, x_A \rangle + \tilde{b}|$ \\ $d(\tilde{w}, B) = |\langle \tilde{w}, x_B \rangle + \tilde{b}|$}

\begin{itemize}
\item Propriétés d'un hyperplan canonique :
\begin{encadre}
\[ m = \frac{1}{\|w\|} \]
\end{encadre}
\couleur{rouge}{\} m est \souligne{fixé} \\ $\hookrightarrow$ y'a qu'une seule solut° donc il est optimal}
\[ \Rightarrow \]
\begin{encadre}
\[ \min_{x_i} |h(x_i)| = 1 \]
\end{encadre}
\couleur{rouge}{Vu que $x_{SV} = h(x_i) = \pm 1 \quad \forall x_i$}
\end{itemize}

\couleur{rouge}{$\hookrightarrow$ Si on dispose de $x_{SV}$, sachant que $h(x_{SV}) = 1$ \\ $\hookrightarrow$ avec $h(\cdot)$, le séparateur canonique}

\couleur{vert}{Si $h(x_{SV}) \neq 1 = a$ \\ alors on récrit $h(x) = \dfrac{1}{a} \tilde{h}(x)$}
\end{document}
